% ============================================================
%  SY02-39  Drinfeld chtoucas 与 Langlands 对应 —— Tata 研究所讲义 · 中文译注本
%  preamble: 版式、定理环境、常用宏
% ============================================================
\documentclass[11pt,a4paper,fontset=windows]{ctexart}

\usepackage[a4paper,margin=2.6cm,top=2.9cm,bottom=2.7cm]{geometry}
\usepackage{amsmath,amscd,amssymb,amsthm,mathtools,mathrsfs,bm}
\usepackage{graphicx}
\usepackage[normalem]{ulem}
\usepackage{array,booktabs,longtable,multirow}
\usepackage{caption}
\usepackage{enumitem}
\usepackage{xcolor}
\usepackage{fancyhdr}
\usepackage{tikz-cd}
\usepackage[colorlinks=true,linkcolor=blue!55!black,citecolor=blue!55!black,
  urlcolor=blue!55!black,bookmarksnumbered=true,
  pdftitle={Drinfeld chtoucas 与 Langlands 对应},
  pdfauthor={Laurent Lafforgue（著）；数译（译注）}]{hyperref}

% ---------------- 定理环境（与原文共用同一计数器，按节编号） ----------------
\theoremstyle{plain}
\newtheorem{theorem}{定理}[section]
\newtheorem{proposition}[theorem]{命题}
\newtheorem{corollary}[theorem]{推论}
\newtheorem{conjecture}[theorem]{猜想}
\newtheorem{lemma}[theorem]{引理}
\theoremstyle{definition}
\newtheorem{definition}[theorem]{定义}
\theoremstyle{remark}
\newtheorem{remark}[theorem]{注}
\newtheorem{example}[theorem]{例}
\newtheorem*{theorem*}{定理}
\newtheorem*{definition*}{定义}
\newtheorem*{proposition*}{命题}
\newtheorem*{corollary*}{推论}
\newtheorem*{lemma*}{引理}
\newtheorem*{remark*}{注}
\newtheorem*{conjecture*}{猜想}
\numberwithin{equation}{section}

% ---------------- 版式 ----------------
\setlength{\parindent}{2em}
\linespread{1.30}
\setlist{nosep,leftmargin=2.2em}
\renewcommand{\arraystretch}{1.15}
\captionsetup{font=small,labelsep=quad}

\pagestyle{fancy}
\fancyhf{}
\fancyhead[L]{\small\kaishu Drinfeld chtoucas 与 Langlands 对应}
\fancyhead[R]{\small\thepage}
\renewcommand{\headrulewidth}{0.4pt}
\fancypagestyle{plain}{\fancyhf{}\fancyhead[R]{\small\thepage}%
  \renewcommand{\headrulewidth}{0pt}}

\ctexset{
  section/format={\Large\bfseries\heiti},
  subsection/format={\large\bfseries\heiti},
  subsubsection/format={\normalsize\bfseries\heiti},
}

\newcommand{\booktitle}{Drinfeld chtoucas 与 Langlands 对应}
\newcommand{\booksubtitle}{Tata 研究所讲义（IHES/M/02/45），附译注}

% 原文印刷页引用
\newcommand{\origpage}[1]{\textsuperscript{\scriptsize\color{gray}[p.~#1]}}

% ---------------- 常用宏 ----------------
\newcommand{\Aa}{\mathbb{A}}
\newcommand{\OO}{\mathcal{O}}
\newcommand{\Fq}{\mathbb{F}_q}
\newcommand{\Fqb}{\overline{\mathbb{F}}_q}
\newcommand{\cE}{\mathcal{E}}
\newcommand{\cF}{\mathcal{F}}
\newcommand{\cG}{\mathcal{G}}
\newcommand{\cH}{\mathcal{H}}
\newcommand{\cL}{\mathcal{L}}
\newcommand{\cM}{\mathcal{M}}
\newcommand{\cN}{\mathcal{N}}
\newcommand{\cO}{\mathcal{O}}
\newcommand{\cP}{\mathcal{P}}
\newcommand{\cX}{\mathcal{X}}
\newcommand{\cZ}{\mathcal{Z}}
\newcommand{\Cht}{\mathrm{Cht}}
\newcommand{\Bun}{\mathrm{Bun}}
\newcommand{\Div}{\mathrm{Div}}
\newcommand{\Loc}{\mathrm{Loc}}
\newcommand{\Rep}{\mathrm{Rep}}
\DeclareMathOperator{\GL}{GL}
\DeclareMathOperator{\PGL}{PGL}
\DeclareMathOperator{\SL}{SL}
\DeclareMathOperator{\Spec}{Spec}
\DeclareMathOperator{\Frob}{Frob}
\DeclareMathOperator{\Gal}{Gal}
\DeclareMathOperator{\Aut}{Aut}
\DeclareMathOperator{\Hom}{Hom}
\DeclareMathOperator{\Pic}{Pic}
\DeclareMathOperator{\rg}{rg}
\DeclareMathOperator{\coker}{coker}
\DeclareMathOperator{\Irr}{Irr}
\DeclareMathOperator{\sgn}{sgn}
\DeclareMathOperator{\Tr}{Tr}
\newcommand{\ZZ}{\mathbb{Z}}
\newcommand{\QQ}{\mathbb{Q}}
\newcommand{\CC}{\mathbb{C}}
\newcommand{\RR}{\mathbb{R}}
\newcommand{\PP}{\mathbb{P}}
\newcommand{\GG}{\mathbb{G}}
\newcommand{\oo}{\infty}
\newcommand{\ul}[1]{\uline{#1}}
\newcommand{\ol}{\overline}
\newcommand{\wh}{\widehat}
\newcommand{\wt}{\widetilde}
\newcommand{\ie}{\emph{i.e.}}
\newcommand{\cf}{\emph{cf.}}
\newcommand{\Frac}{\mathrm{Frac}}

% 中文与长英文词混排时放宽断行，避免 Overfull box
\setlength{\emergencystretch}{3em}
\setcounter{tocdepth}{2}
\renewcommand{\proofname}{证明}

% 原文以罗马数字分章（I--VI），译本沿用
\renewcommand{\thesection}{\Roman{section}}
